\chapter{The XX Chain and Free Fermions}
\label{chap:xx-chain-free-fermions}

The XX chain is the number-conserving limit of the XY mother model.  It is the
cleanest example of how a spin chain can become an ordinary tight-binding
fermion problem after the Jordan--Wigner transformation.  Unlike the TFIM and
the generic XY chain, the XX chain has a conserved fermion number, a Fermi sea,
and a gapless critical phase with two Fermi points.

\section{The \texorpdfstring{$\gamma=0$}{gamma=0} limit of the XY chain}
\label{sec:xx-gamma-zero-limit}

Setting $\gamma=0$ in the transverse-field XY chain gives
\begin{equation}
\label{eq:xx-spin-hamiltonian}
    H_{\rm XX}
    =
    -\frac{J}{2}\sum_{j=1}^{L}
    \left(
        \sigma_j^x\sigma_{j+1}^x
        +
        \sigma_j^y\sigma_{j+1}^y
    \right)
    -h\sum_{j=1}^{L}\sigma_j^z.
\end{equation}
Equivalently, in terms of spin raising and lowering operators,
\begin{equation}
\label{eq:xx-spin-flip-form}
    H_{\rm XX}
    =
    -J\sum_j
    \left(
        \sigma_j^+\sigma_{j+1}^-
        +
        \sigma_j^-\sigma_{j+1}^+
    \right)
    -h\sum_j\sigma_j^z.
\end{equation}
The first term moves a spin flip from one site to a neighboring site.  It does
not create or annihilate spin flips in pairs.  This is the spin-language origin
of fermion-number conservation.

Using the Jordan--Wigner convention of these notes, the bulk fermionic
Hamiltonian is
\begin{equation}
\label{eq:xx-fermion-hamiltonian-real-space}
    H_{\rm XX}^{\rm bulk}
    =
    -J\sum_j
    \left(
        c_j^\dagger c_{j+1}+c_{j+1}^\dagger c_j
    \right)
    -hL+2h\sum_j n_j .
\end{equation}
There is no pairing term.  Therefore
\begin{equation}
\label{eq:xx-number-conservation}
    [H_{\rm XX},N]=0,
    \qquad
    N=\sum_j n_j.
\end{equation}
This is stronger than the parity conservation of the generic XY chain.

For a periodic spin chain, the same parity-sector boundary condition as before
is present:
\begin{equation}
\label{eq:xx-parity-boundary-condition}
    P_F=p
    \quad\Longrightarrow\quad
    c_{L+1}=-p\,c_1.
\end{equation}
In the thermodynamic limit this affects only the quantization of allowed momenta.

\section{Fourier diagonalization}
\label{sec:xx-fourier-diagonalization}

In a fixed boundary sector, use
\begin{equation}
    c_j=\frac{1}{\sqrt{L}}\sum_{k\in\mathcal{K}_\vartheta}
    \ee^{\ii kj}c_k,
\end{equation}
where $\vartheta=0$ or $\pi$ depending on the sector.  The Hamiltonian becomes
\begin{equation}
\label{eq:xx-momentum-hamiltonian}
    H_{\rm XX}
    =
    -hL+
    \sum_{k\in\mathcal{K}_\vartheta}
    \xi(k)c_k^\dagger c_k,
\end{equation}
with single-particle dispersion
\begin{equation}
\label{eq:xx-dispersion}
    \xi(k)=2(h-J\cos k).
\end{equation}
The many-body eigenstates are occupation-number states
\begin{equation}
\label{eq:xx-occupation-basis}
    \ket{\{n_k\}}
    =
    \prod_{k}(c_k^\dagger)^{n_k}\ket{0},
    \qquad
    n_k=0,1,
\end{equation}
with energy
\begin{equation}
\label{eq:xx-many-body-energy}
    E[\{n_k\}]
    =
    -hL+
    \sum_k \xi(k)n_k.
\end{equation}
The diagonalization is therefore simpler than the BdG problem: no Bogoliubov
rotation is required.

\section{Ground-state Fermi sea}
\label{sec:xx-ground-state-fermi-sea}

At zero temperature the ground state fills all modes with negative
single-particle energy:
\begin{equation}
\label{eq:xx-fill-negative-energy}
    \xi(k)<0
    \quad\Longleftrightarrow\quad
    \cos k>\frac{h}{J}
    \qquad (J>0).
\end{equation}
There are three regimes:
\begin{equation}
\label{eq:xx-ground-state-regimes}
\begin{array}{ccl}
    h>J &:& \text{no fermions are occupied},\\[2pt]
    \abs{h}<J &:& \text{a Fermi sea is present},\\[2pt]
    h<-J &:& \text{all fermion modes are occupied}.
\end{array}
\end{equation}
For $\abs{h}<J$, the Fermi points are
\begin{equation}
\label{eq:xx-fermi-momentum}
    k=\pm k_F,
    \qquad
    k_F=\arccos\left(\frac{h}{J}\right).
\end{equation}
The filled region is $\abs{k}<k_F$ in the convention of
\cref{eq:xx-dispersion}.  The fermion density in the thermodynamic limit is
\begin{equation}
\label{eq:xx-fermion-density}
    n
    =
    \langle c_j^\dagger c_j\rangle
    =
    \frac{k_F}{\pi},
    \qquad
    \abs{h}<J.
\end{equation}
Since $\sigma_j^z=1-2n_j$, the magnetization is
\begin{equation}
\label{eq:xx-z-magnetization}
    m_z
    =
    \langle \sigma_j^z\rangle
    =
    1-\frac{2k_F}{\pi},
    \qquad
    k_F=\arccos\left(\frac{h}{J}\right).
\end{equation}
This interpolates from $m_z=-1$ at $h=-J$ to $m_z=+1$ at $h=J$.

The ground-state energy density for $\abs{h}<J$ is
\begin{equation}
\label{eq:xx-ground-state-energy-density}
    e_0
    =
    -h+
    \frac{1}{2\pi}
    \int_{-k_F}^{k_F}2(h-J\cos k)\,\dd k
    =
    -h+
    \frac{2hk_F}{\pi}
    -
    \frac{2J\sin k_F}{\pi}.
\end{equation}
Outside the partially filled regime this expression must be replaced by the
empty or completely filled Fermi sea.

\section{Low-energy Dirac and Luttinger-liquid limit}
\label{sec:xx-low-energy-limit}

In the gapless region $\abs{h}<J$, expand near the two Fermi points:
\begin{equation}
\label{eq:xx-linearization}
    k=\pm k_F+q,
    \qquad
    \abs{q}\ll1.
\end{equation}
The dispersion linearizes as
\begin{equation}
\label{eq:xx-linear-dispersion}
    \xi(+k_F+q)\simeq v_F q,
    \qquad
    \xi(-k_F+q)\simeq -v_F q,
\end{equation}
where
\begin{equation}
\label{eq:xx-fermi-velocity}
    v_F
    =
    \left.\abs{\frac{\dd \xi}{\dd k}}\right|_{k=k_F}
    =
    2J\sin k_F
    =
    2J\sqrt{1-\left(\frac{h}{J}\right)^2}.
\end{equation}
The continuum fermion is decomposed as
\begin{equation}
\label{eq:xx-right-left-decomposition}
    c_j
    \sim
    \ee^{\ii k_F x}\psi_R(x)
    +
    \ee^{-\ii k_F x}\psi_L(x),
    \qquad
    x=ja.
\end{equation}
The low-energy Hamiltonian is the massless Dirac Hamiltonian
\begin{equation}
\label{eq:xx-dirac-continuum-hamiltonian}
    H_{\rm cont}
    =
    -\ii v_F\int\dd x
    \left(
        \psi_R^\dagger\partial_x\psi_R
        -
        \psi_L^\dagger\partial_x\psi_L
    \right).
\end{equation}
In spin-chain language this is the free-fermion point of a Luttinger liquid.  Its
central charge is
\begin{equation}
\label{eq:xx-central-charge}
    c=1,
\end{equation}
which distinguishes the XX critical region from the Ising critical point of the
TFIM, where $c=1/2$.

\section{Correlation functions}
\label{sec:xx-correlation-functions}

Because the Hamiltonian is number conserving and diagonal in momentum space, the
basic equal-time fermion correlator is elementary.  In the thermodynamic ground
state with $\abs{h}<J$,
\begin{equation}
\label{eq:xx-fermion-two-point}
    \langle c_i^\dagger c_j\rangle
    =
    \frac{1}{2\pi}\int_{-k_F}^{k_F}
    \ee^{\ii k(j-i)}\,\dd k
    =
    \begin{cases}
        \dfrac{k_F}{\pi}, & i=j,\\[8pt]
        \dfrac{\sin[k_F(j-i)]}{\pi(j-i)}, & i\neq j.
    \end{cases}
\end{equation}
This $1/r$ decay is the fermionic signature of gaplessness.

The longitudinal spin correlator is local in fermion density:
\begin{equation}
\label{eq:xx-longitudinal-correlator-density}
    \sigma_j^z=1-2n_j.
\end{equation}
The transverse spin correlators, such as
$\langle\sigma_i^x\sigma_j^x\rangle$, contain Jordan--Wigner strings and are
computed as determinants or Pfaffians.  Their long-distance behavior in the
critical region is algebraic, rather than exponential.  At half filling
($h=0$, $k_F=\pi/2$), the leading transverse spin correlations decay as a power
law with the characteristic free-fermion/Luttinger-liquid exponent.

\section{Comparison with the generic XY chain}
\label{sec:xx-comparison-generic-xy}

The XX chain is obtained from the generic XY chain by turning off the pairing
amplitude.  The distinction can be summarized as follows.

\begin{table}[h]
\centering
\small
\setlength{\tabcolsep}{3pt}
\caption{Generic XY chain versus the XX limit.}
\label{tab:xx-vs-xy}
\begin{tabularx}{\textwidth}{L{0.25\textwidth}L{0.33\textwidth}Y}
\toprule
Feature & Generic XY chain $(\gamma\neq0)$ & XX chain $(\gamma=0)$ \\
\midrule
Fermion number & Not conserved & Conserved \\
Fermion parity & Conserved & Conserved, with full $U(1)$ enhanced symmetry \\
Momentum-space block & BdG $(k,-k)$ block & Single-particle dispersion $\xi(k)$ \\
Ground state & BCS paired state & Fermi sea \\
Generic criticality & Ising lines at $h=\pm J$ & Gapless region $\abs{h}<J$ \\
Continuum theory & Majorana, $c=1/2$ & Dirac/Luttinger liquid, $c=1$ \\
\bottomrule
\end{tabularx}
\end{table}

Adding a small nonzero anisotropy $\gamma$ to the gapless XX region introduces a
$p$-wave pairing perturbation.  In the fermionic language this pairs the right-
and left-moving modes and opens a gap except where the lattice mass also closes
at the Ising transition.

\section{Summary}
\label{sec:xx-summary}

The XX chain is the $\gamma=0$ limit of the XY mother model.  It is exactly a
number-conserving tight-binding model with dispersion
\begin{equation}
    \xi(k)=2(h-J\cos k).
\end{equation}
For $J>0$, the ground state is empty for $h>J$, completely filled for $h<-J$,
and a partially filled Fermi sea for $\abs{h}<J$.  The latter regime has two
Fermi points, a linear low-energy Dirac description, and central charge $c=1$.
This chapter therefore illustrates the difference between free-fermion
solvability by a Fermi sea and free-fermion solvability by BdG pairing.
